\documentclass[11pt]{article}
\usepackage[margin=25mm]{geometry}
\usepackage{amsmath,amssymb,hyperref}
\title{Nested vertex and disjoint primitive: exact reduction and an integrable rotation control}
\author{Phi3Lab development record}\date{25 October 2026}
\begin{document}\maketitle
\noindent This derivation applies to G56, G57, G201, G210, G215, G71 and G252. It is not a numerical quadrature error certificate. The mass is one; each six-dimensional loop has measure $d^6q/\pi^3$. Couplings and graph symmetry factors are excluded.
\paragraph{Retaining the parent subtraction.}
Two primitive loops are integrated analytically after their local subtractions. With $H_o$ the three outer propagators and $H_i$ the three parent propagators, the full integrand is
\[
(1-T^0_{\rm phys})H_o O\left[H_iT_i-K_0(k)\right],\qquad
K_0(k)=-\frac{E(k^2)}{2(1+k^2)^3},\quad E(s)=\int_0^1\log(1+sx(1-x))\,dx.
\]
Here $T_i=-L(a,b,c)$ is the complete primitive triangle and $O$ is either another triangle or $J(s)=B(s)/(1+s)$. In the bubble case exactly one matching host denominator enters $J$; the independent reference keeps all original host denominators and $B$.
The parent's subtraction sets its boundary momenta to zero but retains $k$. The physical subtraction sets only $p_1,p_2$ to zero and retains $q,k$. The parent branches must stay combined.
\paragraph{Original integrability.}
All $2^{12}-1=4095$ nonempty internal-edge subsets of each original graph are enumerated independently. For every connected 1PI loop subgraph,
\[
\ell=I-V+1,\qquad n=3V-2I,\qquad \omega=6\ell-2I=6-2n.
\]
All proper subgraphs with $\omega\ge0$ and their Taylor degrees match the original recursive program. There are eight forests per graph. Gaussian, positive-polynomial and Schur constructions agree exactly on all 56 states. These massive Euclidean graphs with zero-momentum, inner-first BPH subtractions satisfy the hypotheses of the massive BPH convergence theorem. The theorem is an input; the scripts check its graph-specific hypotheses, not a new proof of that theorem. Integration of the subtracted primitive loops preserves absolute integrability.
\paragraph{Damped quadratic control.}
Let $C$ denote the coefficient of $\lambda^2$ in the full parent-subtracted product under $p\mapsto\lambda p$. Subtract $\chi C$ and return its exact rotational mean, with the predeclared invariant
\[
\chi(q,k)=\left[1+\left(\frac{q^2+k^2}{64}\right)^4\right]^{-3}.
\]
This changes neither local nor overall Taylor subtraction. Put $R=1+q^2+k^2$, $L=\log(3R)$. The three Taylor coefficients of one baseline-normalized propagator have bounds $(1,3\sqrt R,9R)$. For six propagators these become $(1,18\sqrt R,189R)$. Each primitive factor is bounded by $(64L,64\sqrt R L,64RL)$. Since $|K_0|\le L$,
\[
|C|\le (933888+4096)RL^2=937984RL^2<2^{20}RL^2.
\]
This bound is used only for the damped control. It does not separate the undamped divergent parent branches.
Writing $T=q^2+k^2$, $x=q^2/T$ and integrating $x$ gives the radial measure $T^5dT/120$. For $T\ge64$ the control is dominated by
\[
\frac{2^{21}64^{12}}{120}\,T^{-6}\log^2(6T),\qquad
\int_a^\infty\frac{\log^2(6T)}{T^6}\,dT
=\frac1{a^5}\left(\frac{\log^2(6a)}5+\frac{2\log(6a)}{25}+\frac2{125}\right).
\]
The compact region is finite. Thus both the added control and its rotation average are absolutely integrable.
\paragraph{Independent rotation check.}
The loop Gram matrix is retained:
\[
\left\langle(q\cdot p_a)(k\cdot p_b)\right\rangle
=\frac{(q\cdot k)(p_a\cdot p_b)}6.
\]
Native code uses this covariance. The independent reference uses formal Taylor convolution and 24 explicitly constructed six-dimensional frames. Three unimodular charts cover the $q,k,q+k$ soft regions. An optional inner-loop reflection leaves physical momenta untouched.
\paragraph{Reproducibility.}
The exact proof is in \texttt{verify\_nested\_disjoint\_control\_v1.py}; its durable receipt records 56 forest states and 28665 subsets. Native tests, diagnostics, production, independent replay and accounting have separate exit records. No precision acceptance follows from this derivation alone.
\paragraph{Mathematical reference.}
H. van Hees, \emph{Introduction to Relativistic Quantum Field Theory}, sections 5.7--5.9, equations (5.215)--(5.222):\newline\url{https://itp.uni-frankfurt.de/~hees/publ/lect.pdf}.
\end{document}
